\section{Barycentric allocation and the deployed law}\label{sec:allocation}
We give the finite-target form of the classical convex-order argument explicitly, since it is also the building block for phase-compatible implementation.

\begin{lemma}[Finite barycentric allocation]\label{lem:allocation}
Let $K$ be compact metrizable and convex as above, $\lambda\in\mathcal P(K)$, and $\nu=\sum_{j=1}^m p_j\delta_{q_j}$. The following are equivalent:
\begin{enumerate}
\item there are Borel functions $r_j:K\to[0,1]$, $\sum_jr_j=1$ $\lambda$-almost everywhere, such that
\begin{equation}\label{eq:moments}
 \int r_j\,d\lambda=p_j,\qquad
 \int r_j(w)f(w)\lambda(dw)=p_jf(q_j)\quad(f\in\aff(K));
\end{equation}
\item all inequalities \eqref{eq:localdual} hold;
\item $\nu\cx\lambda$.
\end{enumerate}
When the inequalities fail, finitely many determining affine coordinates and $m$ affine expressions provide a strict witness. Atomic acquired laws are allowed, and randomization is not silently replaced by deterministic cells.
\end{lemma}
\begin{proof}
Consider the simplex of allocations $r=(r_1,\ldots,r_m)$ in $(L^\infty(\lambda))^m$. It is weak-star compact: it is contained in the unit ball, and positivity and the identity $\sum_jr_j=1$ are weak-star closed. Map it to all its masses and affine-coordinate moments. Each coordinate is an integral against an $L^1(\lambda)$ function, so the image is compact and convex in a product of real lines. If the target moments $(p_j,p_jq_j)$ are outside that image, separation in this locally convex product gives a continuous linear functional involving only finitely many coordinates. Combining the coordinates for each $j$ yields an affine $\phi_j$. The supremum of the separated functional over allocations is exactly
\[
 \int\max_j\phi_j(w)\lambda(dw):
\]
a measurable least maximizing index attains it. This proves equivalence of the first two assertions, including strict separation. A weak-star allocation has Borel versions on the standard Borel space; modify on one null set to keep it in the simplex everywhere. Matching the separating coordinates matches the barycenter, hence every continuous affine function.

Given \eqref{eq:moments}, generate $W\sim\lambda$ and choose label $j$ with probability $r_j(W)$. Then $\E[W\mid j]=q_j$, where conditional expectation is understood through affine coordinates. Jensen's inequality gives $\nu\cx\lambda$. Conversely, for any affine family, convex order gives
\[
 \sum_jp_j\phi_j(q_j)\le\int\max_i\phi_i\,d\nu
 \le\int\max_i\phi_i\,d\lambda,
\]
which is assertion two. This also proves sufficiency without invoking a separate infinite-dimensional transport theorem. The argument is the finite-target martingale/convex-order mechanism of Strassen\cite{strassen}, not a new static comparison theorem.
\end{proof}

\begin{lemma}[Calibration under all incoming reports]\label{lem:calibration}
Suppose label $i$ has probability $p_i$ and conditional predictive state $q_i$. Draw its action from $\alpha_i$, execute the raw report kernel, and let $W=F_t(q_i,a,Y)$. If a next-label rule satisfies $\E[W\mid Z'=j]=q'_j$, then $q'_j$ is the actual predictive state conditional on $Z'=j$. In particular this applies to the kernel $r_j(W)$ of Lemma~\ref{lem:allocation} under the full acquired mixture $\lambda=\sum_{i,a}p_i\alpha_i(a)\Lambda_t(q_i,a)$.
\end{lemma}
\begin{proof}
The incoming label represents a conditional mixture of physical histories, not a physical preparation performed anew. Equation~\eqref{eq:affine_instrument} and predictive sufficiency imply that the conditional expectation of each determining future test, given $(i,a,Y)$, is its affine evaluation at $W$. Integrating this identity over \emph{all} incoming labels and reports sent to $j$ gives its evaluation at $q'_j$. The tests determine the predictive quotient. This is precisely the conditional state under the deployed observer. There is no comparison with the full-history posterior on a different controlled path.
\end{proof}

\begin{proof}[Proof of Theorem~\ref{thm:completion}(a), (c) for the external interface, and (d)]
For necessity, at each persistent cut of any observer take its actual conditional predictive states. Their weighted law is $\nu_t$, with support size at most $|Z_t|$. If several labels have the same state, average their action distributions using their conditional label probabilities. This gives the same acquired mixture $\lambda_{t+1}$. The next retained state is a conditional barycenter of the pre-pooling state, so Jensen's inequality gives \eqref{eq:flow}. This step records only marginals needed for subsequent predictive evolution; it does not assert equality of the original labelled trajectories after duplicate states are merged.

For sufficiency, assume a flow. Use its atoms as phase labels and its action probabilities as the programmed action kernel. Lemma~\ref{lem:allocation} provides $r_{t+1,j}$ for $\lambda_{t+1}\to\nu_{t+1}$. Programme
\[
 k_t(j\mid i,a,y)=r_{t+1,j}(F_t(q_{t,i},a,y)).
\]
The conditional-state parameters are constants in this phase's programme, not a persistently updated real vector. Lemma~\ref{lem:calibration} proves the prescribed occupancies and states by induction from $q_0$. Thus each feasible flow yields a legal observer of risk $\int g\,d\nu_n$. Conversely, replacing any observer's final readout by the best conditional readout gives this risk. Infimizing proves the first equality in \eqref{eq:exactrisk}.

For the telescope, let $Q_t$ denote the actual retained conditional state. Before pooling at the next cut,
\[
 \E V_{t+1}(W_{t+1})-\E V_t(Q_t)=A_t\ge0
\]
by the Bellman minimum. Calibration and concavity give
\[
 \E V_{t+1}(Q_{t+1})-\E V_{t+1}(W_{t+1})=J_t\ge0.
\]
Adding and telescoping yields \eqref{eq:slack}. The argument uses the machine's actual action distribution and remains valid at nonsmooth action switches. It neither compares independently optimized shadow controllers nor assumes Bellman-optimal actions for a stationary machine that reuses a label.
\end{proof}

\begin{remark}[Why the dual is operational]\label{rem:dual}
An arbitrary terminal splitting of a preparation is not automatically obtainable from a fixed sensor. The right side of \eqref{eq:localdual} is evaluated under that sensor's actual $\lambda$, not under an unconstrained information-design experiment. Likewise, local contractions at different cuts do not authorize phase-dependent code. The simultaneous test below imposes precisely the shared-kernel restriction that ordinary one-step convex order omits.
\end{remark}

\subsection{Finite obstructions to a global risk target}
For the externally clocked problem, compactness upgrades the witness for one proposed flow to a finite obstruction valid for every flow at an unattainable risk level. This is not an interchange of an infimum with a pointwise dual supremum.

\begin{theorem}[Compact completion and finite obstructions]\label{thm:finite_obstruction}
Under the compactness and weak-continuity hypotheses of Theorem~\ref{thm:completion}, every finite external width profile has an optimal allocation flow and an optimal measurable observer. For each phase fix a countable family $h_{t,1},h_{t,2},\ldots$ dense in the continuous convex functions on $K_t$, each a finite maximum of finite affine combinations of the determining coordinates. Let $c_k$ be the minimum terminal risk subject to the atomic-width, action-marginal and acquisition constraints of \eqref{eq:flow}, but retaining only the first $k$ convex-order inequalities at each cut. Then
\begin{equation}\label{eq:finite_obstruction}
 c_k\uparrow R_{\ext}(\mathbf m).
\end{equation}
Consequently, if a number $r<R_{\ext}(\mathbf m)$ is proposed as a feasible risk, finitely many of these acquired-law inequalities, together with the resource and risk constraints, already exclude every proposed flow with risk at most $r$. This is a finite family of tests, not necessarily a single test that defeats every candidate. No effective bound on its size or computability is asserted for general compact Borel data.
\end{theorem}
\begin{proof}
Represent the action assignments by probabilities $\gamma_t$ on $K_t\times\mathcal A_t$ whose first marginal is $\nu_t$. The sets of probabilities with at most $m_t$ atoms are compact: a sequence of such probabilities has convergent subsequences of weights and atom locations after padding by zero weights. Marginal constraints are closed. The map
$\gamma_t\mapsto\lambda_{t+1}=\int\Lambda_t(q,a)\gamma_t(dq,da)$
is weakly continuous by weak continuity of each $\Lambda_t$ and finiteness of the action set. Thus the candidate space before imposing convex order is compact. Every convex-order inequality is closed, and terminal risk is continuous. Feasibility is nonempty: choose one legal action at every cut and replace each acquired law by its single barycenter. Hence a risk minimizer exists and is synthesized by Theorem~\ref{thm:completion}.

For completeness, the required countable polyhedral family exists. The closed convex epigraph of a continuous convex function on compact $K_t$ can be separated from each point lying strictly below it by a continuous affine functional of the ambient countable product and the height coordinate. The height coefficient cannot vanish, because the epigraph contains a point with the same base coordinate. The resulting affine minorant uses only finitely many coordinates. A finite subcover of $K_t$ gives a finite maximum approximating the function uniformly from below. Approximating its finitely many coefficients by rationals gives a countable uniformly dense collection, with constants included. Inequalities for this collection imply all continuous convex inequalities by uniform approximation.

The finite-constraint feasible sets are nonempty, compact and nested, so their minima $c_k$ are nondecreasing and at most the full optimum. Choose minimizers and a convergent subsequence in the compact candidate space. For every fixed test index, all sufficiently late members obey its inequality, and the limit does too. Thus the limit satisfies all convex-order constraints, while its terminal risk is $\lim_kc_k$. This proves \eqref{eq:finite_obstruction}. If $r$ is strictly below that limit, some finite $k$ has $c_k>r$, yielding the asserted obstruction. Its proof does not bound $k$ and does not establish a finite computational representation of the carrier or kernels.
\end{proof}
